\section{Discussion}
\label{sec:discussion}

Hiddenness is a global property, and training controls a model only near its data. An attractor stays hidden in a model
only if the local unstable manifold of no unstable equilibrium of the model reaches its basin (Corollary~\ref{cor:criterion}).
The model is fitted on the attractor and, at best, on transients toward the visible state; nothing in the loss constrains
where the unstable manifolds of equilibria away from the data go. Topology fixes the bookkeeping of these equilibria: their
number is odd, every additional pair contains an unstable one (Corollary~\ref{cor:counts}), and when the true system has
$s\neq1$, reproducing its equilibria forces additional ones (Proposition~\ref{prop:forced}). How many pairs a model carries is
set by training, and the measured survival curves show how quickly the resulting unstable equilibria cost hiddenness: one in
the region costs about two thirds of the probability of staying hidden, and with nine or more one to three reproduced
attractors in a hundred survive. The decline is similar for echo state networks, one- and two-layer perceptrons and residual networks once
the equilibria are counted where the dynamics lives, which identifies $\Uin$ as the main variable; the class-specific slopes
and the effect of the training set show that it is not the only one.

The two cases with $s\neq1$ show what topology does and does not decide. In Chua's circuit, regime~I,
Proposition~\ref{prop:forced} does not forbid a bounded model to copy all three true equilibria; it requires additional ones
if it does. The bounded models in our ensembles instead miss one of the saddle-foci, while residual networks copy all three
in $23$ of $480$ cases; but the fixed points of a residual network are the zeros of $g\approx h f$, which biases it toward the
equilibria of the flow, so this comparison mixes topology with architecture. Where bounded models have extra equilibria,
these are mostly stable and do not by themselves destroy hiddenness. In the equilibrium-free Li--Sprott system every
bounded model has an equilibrium by Brouwer's theorem, which the true system does not have, and in every one of the $29$
reproductions the training attractor is excited from it. Topology forces the equilibrium, not its excitation; but in all our reproductions the attractor, which has
no equilibrium to be excited from in the true system, is self-excited in the model, an aspect of the reconstruction of
unseen attractors\cite{Rohm2021} that the reconstruction of the attractor itself does not reveal.

Several practical rules follow. After training, the equilibria of a model should be located and, for bounded models, their
index sum checked; Corollary~\ref{cor:counts}(iii) turns this into a free test of the root finding. The number of unstable
equilibria in the region of the true system is a cheap predictor of whether hiddenness survives. With $\Uin=0$ only equilibria outside the
region can excite the attractor, and $94\%$ of the reproductions by perceptrons and residual networks stayed hidden; each
further unstable equilibrium lowers the chance, from $29\%$ at $\Uin=1$ to $13\%$ at $\Uin=3$--$4$ and $1\%$ at nine or more. Survival is not fidelity, however: a model without the visible
equilibrium, or with a basin many times too large, can keep the attractor hidden for the wrong reason, so hiddenness should
be checked together with the basins and the visible state (Sec.~\ref{sec:fidelity}). Selecting among trained
models by $\Uin$ is therefore more useful for this purpose than increasing the size of the network, which acts mainly through
$U$ and can increase it. Removing unstable equilibria during training also keeps the attractor hidden: a penalty on their
spectral radius did so in $0.994$ of the reproductions at an unchanged median one-step error (Sec.~\ref{sec:penalty}). This
optimizes hiddenness rather than fidelity, however: by Corollary~\ref{cor:criterion} a model without unstable equilibria
in the region can only be excited from outside it, and the penalty removes every unstable equilibrium in the region, true
ones included, such as the saddle-foci of Chua's circuit and the unstable half-line of LE1. It is a prior that suits systems
known to have no unstable equilibria there; otherwise the known equilibria must be exempted from it.
Whether a given unstable equilibrium excites the attractor is decided by its unstable manifold, so the excitation test can
sample the unstable subspace instead of a ball in the full state space; for high-dimensional models this saves more than an
order of magnitude of trajectories, at the price of missing basins reached from a tiny fraction of the directions. Including
transients toward the visible attractor in the training set raises survival, in echo state networks for one half to three
quarters through fewer unstable equilibria and in the model classes also at fixed $\Uin$. Finally, a model should not be trusted to reveal a hidden attractor that its data do
not approach, and imposing the true equilibria on the model does not make it more faithful: in our echo state networks it
added unstable equilibria.

These results give a precise form to several observations in the literature. Untrained attractors have been argued to be
intrinsic to learning systems with a bounded state space;\cite{OHagan2025Confabulation} Theorem~\ref{thm:degree} and
Proposition~\ref{prop:forced} make this exact for equilibria and identify when they are forced. Spurious slow modes that
pull closed-loop reservoirs to fixed points\cite{Oishi2026FastSlow} and extra saddles organizing a reservoir's
attractor\cite{Kabayama2025PseudoLorenz} are instances of the spurious equilibria counted here, and the dependence of basin
predictions on warm-up and model structure\cite{ZhangCornelius2023} reflects the same freedom of a model away from its data.
Architectures that guarantee bounded or dissipative learned
dynamics\cite{Kaptanoglu2021Trapping,Linot2023Stabilized,Tang2026Dissipative} remove escape at the cost described by
Proposition~\ref{prop:forced}. The closed-loop echo state network $r\mapsto\tanh(Ar+b)$ is a discrete Lur'e system, the class
in which hidden oscillations were first studied through the Aizerman and Kalman
problems;\cite{LeonovKuznetsov2013,Kuznetsov2019KalmanHidden} trained networks can also have hidden attractors of their own,
which are spurious, as hand-designed Hopfield networks do,\cite{DancaKuznetsov2017Hopfield} and we leave them to a separate
study.

The study has limits. All true systems are low-dimensional, although the models live in up to $2000$ dimensions. The large
ensembles are labeled by the light protocol, whose ``not excited'' marks a candidate; it agreed with the full protocol in all $115$
pairs of the six regimes, and on the benchmark the full protocol itself mislabeled one of $46$ attractors as hidden. A missed exciting
equilibrium would lower $U$ and make the attractor look hidden at the same time; the index-sum check excludes incomplete sets
of bounded models, but a missed pair of opposite indices passes it. Searches with up to eight times the starts found
further equilibria only in two of $184$ echo state networks, both with $\rho=1.2$ and without a reproduced attractor, and
the certified enumeration found none in $960$ three-dimensional models; the largest reservoirs, however, are represented by only eight networks each for $N=1400$
and $2000$. For echo state networks hiddenness is defined in the
reservoir space, the state space of the model as a dynamical system; a user who initializes the network by synchronization
reaches only a neighborhood of the synchronization manifold, the image of the echo state map, which is generically an
embedding of the driving system,\cite{HartHookDawes2020} and we have not measured how far the equilibria of the networks
lie from it. The perceptrons, which act in the physical space, show a similar decline. The survival curve is
an empirical regularity with a heuristic interpretation, not a theorem: equilibria excite the attractor more often than
points at the same distance lie in its basin in some regimes and not in others, their excitations are correlated within a
model, and the slope of the curve differs between regimes by a factor of about two. Residual networks lie outside
Theorem~\ref{thm:degree}, and counting their equilibria in the region of the true system is a choice justified by the data.
The penalty experiment used one architecture, two-layer perceptrons of width $256$, and moved almost all models to
$\Uin=0$: it tests the step from $\Uin\ge1$ to $\Uin=0$ rather than the curve at $\Uin\ge1$, and a variant that spares the
known true equilibria was not run.
